% cpu/swdesign.tex
% mainfile: ../perfbook.tex
% SPDX-License-Identifier: CC-BY-SA-3.0

\section{软件设计启示}
\label{sec:cpu:Software Design Implications}
%
\epigraph{一船向东行，一船向西驶，\\
	  吹的却是同样的风；\\
	  决定航向的不是风的力量，\\
	  而是帆的方向。}
	 {Ella Wheeler Wilcox}

\cref{tab:cpu:CPU 0 View of Synchronization Mechanisms on 8-Socket System With Intel Xeon Platinum 8176 CPUs at 2.10GHz}
中比率的数值至关重要，因为它们限制了并行应用程序的效率。
为了弄清这一点，我们假设一款并行程序使用\IXacr{cas}操作来进行线程间通信。
假设进行通信的线程不是与自身而是主要与其他线程通信，
那么CAS操作通常会涉及cache miss。
进一步假设对应每次CAS通信操作的工作需要消耗300\,ns，
这足以计算几个浮点超越函数。
其中一半的执行时间将消耗在CAS通信操作上，
这也意味着两个CPU的系统运行这样一个并行程序的速度，
并不比单CPU系统运行一个串行程序的速度快。

在分布式系统的情况下结果还会更糟，
因为单次通信操作的延迟可能和数千甚至数百万次
浮点运算所需的时间一样长。
这也说明了通信操作应该尽力避免，每次通信应该能够处理大量数据。

\QuickQuiz{
	既然分布式系统的通信如此昂贵，
	为什么还有人费心使用这样的系统？
}\QuickQuizAnswer{
	原因有很多：

	\begin{enumerate}
	\item	共享内存多处理器系统有严格的规模限制。
		如果你需要超过几千个CPU，就别无选择，
		只能使用分布式系统。
	\item	大型共享内存系统的每单位计算成本
		往往高于小型系统。
	\item	大型共享内存系统的cache miss延迟
		往往比小型系统大得多。
		要看到这一点，请比较
		\cref{tab:cpu:CPU 0 View of Synchronization Mechanisms on 8-Socket System With Intel Xeon Platinum 8176 CPUs at 2.10GHz}
		（第~\cpageref{tab:cpu:CPU 0 View of Synchronization Mechanisms on 8-Socket System With Intel Xeon Platinum 8176 CPUs at 2.10GHz}页）
		和\cref{tab:cpu:CPU 0 View of Synchronization Mechanisms on 12-CPU Intel Core i7-8750H CPU @ 2.20GHz}。
	\item	分布式系统的通信操作不一定消耗太多CPU，
		因此计算可以与消息传输并行进行。
	\item	许多重要问题是``令人尴尬的并行''，
		因此极少量的消息就能启用极大量的处理。
		SETI@HOME~\cite{SETIatHOME2008}
		就是此类应用的一个例子。
		这类应用即使在通信延迟极长的情况下，
		也能很好地利用计算机网络。
	\end{enumerate}

	因此，大型共享内存系统往往用于那些受益于比分布式计算
	所能提供的更快延迟的应用，
	特别是那些受益于大共享内存的应用。

	继续进行的并行应用研究工作很可能会增加能够在
	通信延迟较长的机器和/或集群上良好运行的令人尴尬的并行
	应用的数量，降低成本是推动这一趋势的驱动力。
	话虽如此，大幅降低硬件延迟将是一个极受欢迎的发展，
	无论是对于单系统还是分布式计算。
	事实上，这样的降低正在发生，
	尽管从我们这些经历过2000年代初期之前
	CPU时钟频率指数增长的人的角度来看，速度相当缓慢。
}\QuickQuizEnd

这一课应该非常明了：
并行算法必须将每个线程设计成尽可能独立运行的线程。
越少使用线程间通信手段——比如\IX{atomic}操作、lock或者其他消息传递方法——
应用程序的性能和可扩展性就越好。
这种方法将在\cref{chp:Counting}中初步涉及，
在\cref{chp:Partitioning and Synchronization Design}中深入探讨，
并在\cref{chp:Data Ownership}中推向其逻辑极致。

另一种方法是确保任何共享都是读多写少的，
这允许CPU的cache复制读多写少的数据，
从而使所有CPU都能快速访问。
这种方法在\cref{sec:count:Eventually Consistent Implementation}中初步涉及，
并在\cref{chp:Deferred Processing}中更深入地探讨。

简而言之，想要达到优秀的并行性能和可扩展性，
就意味着在并行算法和实现中努力追求，
无论是通过精心选择数据结构和算法、利用现有的并行软件和环境，
还是将并行问题转化为已经有并行解决方案存在的问题。

\QuickQuiz{
	好吧，如果我们不得不将分布式编程技术
	应用于共享内存并行程序，
	那为什么不总是使用这些分布式技术而抛弃共享内存呢？
}\QuickQuizAnswer{
	因为通常情况下，只有一小部分程序是性能关键的。
	共享内存并行性允许我们将分布式编程技术
	集中在那一小部分上，
	而在非性能关键的大部分程序上使用更简单的共享内存技术。
}\QuickQuizEnd
